% Part: sets-functions-relations
% Chapter: functions
% Section: basics

\documentclass[../../../include/open-logic-section]{subfiles}

\begin{document}

\olfileid{sfr}{fun}{bas}
\olsection{मूलभूत संकल्पना}

\begin{explain}
दिलेल्या संचातील प्रत्येक !!{element} दुसऱ्या एखाद्या दिलेल्या संचातील
एका ठरावीक !!{element} शी जोडणारी संगती म्हणजे \emph{फलन}.
उदाहरणार्थ, $1$ मिळवण्याच्या क्रियेने एक फलन ठरते: प्रत्येक संख्या~$n$
एका एकमेव संख्या~$n+1$ शी जोडली जाते.

अधिक व्यापकपणे पाहता, फलनांना आदान म्हणून जोड्या, त्रिके इत्यादी देता
येतात आणि त्यांच्याकडून कोणत्या तरी प्रकारचे प्रदान मिळते. प्राथमिक
अंकगणितातून अनेक फलने आपल्याला परिचित आहेत. उदाहरणार्थ, बेरीज आणि
गुणाकार ही फलने आहेत. ती दोन संख्या आदान म्हणून घेतात आणि तिसरी संख्या देतात.

या गणिती, अमूर्त अर्थाने फलन म्हणजे एक \emph{काळी पेटी}: कोणत्या आदानाशी
कोणते प्रदान जोडलेले आहे एवढेच महत्त्वाचे असते; प्रदानाची गणना करण्याची
पद्धत महत्त्वाची नसते.
\end{explain}

\begin{defn}[फलन]
\emph{फलन} $f \colon A \to B$ ही~$A$ मधील प्रत्येक !!{element} ची~$B$
मधील एका !!{element} शी लावलेली संगती असते.

$A$ ला~$f$ चा \emph{प्रांत} आणि $B$ ला~$f$ चा \emph{सहप्रांत} म्हणतो.
~$A$ मधील !!{element}s ना~$f$ ची आदाने किंवा \emph{फलसाधक} म्हणतात.
फलसाधक~$x$ शी~$f$ ने जोडलेल्या~$B$ मधील !!{element} ला फलसाधक~$x$
साठीचे \emph{$f$ चे मूल्य} म्हणतात आणि ते~$f(x)$ असे लिहितात.

~$f$ ची \emph{व्याप्ती} $\ran{f}$ म्हणजे सहप्रांताचा असा उपसंच की
त्याचे घटक कोणत्या तरी फलसाधकासाठीची~$f$ ची मूल्ये असतात;
$\ran{f} = \Setabs{f(x)}{x \in A}$.
\end{defn}

फलनांचा विचार करण्यासाठी \olref{fig:function} मधील आकृती उपयोगी पडू शकते.
डावीकडील लंबवर्तुळ फलनाचा \emph{प्रांत} दर्शवते; उजवीकडील लंबवर्तुळ
त्याचा \emph{सहप्रांत} दर्शवते; आणि बाण प्रांतातील \emph{फलसाधकापासून}
सहप्रांतातील त्याच्या संगत \emph{मूल्याकडे} जातो.

\begin{figure}
  \olasset{assets/diagrams/function.tikz}
  \caption{फलन एका संचातील प्रत्येक !!{element} ची दुसऱ्या संचातील
    !!a{element} शी संगती लावते. बाण प्रांतातील फलसाधकापासून
    सहप्रांतातील त्याच्या संगत मूल्याकडे जातो.}
  \ollabel{fig:function}
\end{figure}

\begin{ex}
गुणाकार नैसर्गिक संख्यांच्या जोड्या आदान म्हणून घेतो आणि त्यांना प्रदान
म्हणून नैसर्गिक संख्या जोडतो. म्हणून तो $\Nat \times \Nat$ या प्रांताकडून
$\Nat$ या सहप्रांताकडे जातो. त्याची व्याप्तीही $\Nat$ च असते, कारण
प्रत्येक $n \in \Nat$ हा $n \times 1$ असतो.
\end{ex}

\begin{ex}
गुणाकार हे फलन आहे, कारण तो प्रत्येक आदानाशी---नैसर्गिक संख्यांच्या
प्रत्येक जोडीशी---एकच प्रदान जोडतो: $\times \colon \Nat^2 \to \Nat$.
याउलट, नैसर्गिक संख्या~$n$ शी $x^2 = n$ असणारी वास्तव संख्या~$x$
जोडण्याची संगती फलन होत नाही, कारण प्रत्येक धन पूर्णांक~$n$ ची दोन
वर्गमुळे असतात: $\sqrt{n}$ आणि~$-\sqrt{n}$. फक्त धन वर्गमूळ देऊन
आपण तिचे फलन करू शकतो: फक्त ऋणेतर (प्रधान) वर्गमूळ परत द्यायचे,
$\sqrt{\phantom{X}} \colon \Nat \to \Real$.
\end{ex}
%READERNOTE{OLFUN-002 — गोठवलेल्या इंग्रजी स्रोतामध्ये सर्व नैसर्गिक संख्यांसाठी positive square root म्हटले आहे. शून्याचे प्रधान वर्गमूळ शून्य असल्यामुळे येथे ऋणेतर (प्रधान) असे दुरुस्त केले आहे; धन पूर्णांकांना दोन वास्तव वर्गमुळे असतात हे आधीचे विधान योग्य आहे.}

\begin{ex}
वर्गातील प्रत्येक विद्यार्थ्याशी त्याची अंतिम श्रेणी जोडणारा संबंध हे
फलन आहे---एकाच वर्गात कोणत्याही विद्यार्थ्याला दोन वेगवेगळ्या अंतिम
श्रेणी मिळू शकत नाहीत. वर्गातील प्रत्येक विद्यार्थ्याशी त्याचे पालक
जोडणारा संबंध फलन नाही: विद्यार्थ्यांचे पालक शून्य, दोन किंवा अधिक असू शकतात.
\end{ex}

\begin{explain}
प्रत्येक शक्य फलसाधकासाठी फलनाचे मूल्य काय असते हे नेमक्या पद्धतीने
सांगून फलनाची व्याख्या करता येते. सूत्र देणे, मूल्याची गणना करण्याची
पद्धत वर्णन करणे किंवा प्रत्येक फलसाधकासाठीची मूल्ये यादीत देणे या
त्याच्या वेगवेगळ्या पद्धती आहेत. फलनाची व्याख्या कोणत्याही पद्धतीने
केली, तरी प्रत्येक फलसाधकासाठी एक आणि केवळ एकच मूल्य दिलेले असेल
याची खात्री केली पाहिजे.
\end{explain}


\begin{ex}
$f \colon \Nat \to \Nat$ ची व्याख्या $f(x) = x+1$ अशी करू.
या व्याख्येत $f$ हे नैसर्गिक संख्या आदान म्हणून घेऊन नैसर्गिक संख्या
प्रदान करणारे फलन आहे असे सांगितले आहे. नैसर्गिक संख्या~$x$ दिली,
की $f$ तिची उत्तरवर्ती संख्या~$x+1$ देईल असे ती सांगते.
या उदाहरणात सहप्रांत $\Nat$ हा~$f$ ची व्याप्ती नाही, कारण नैसर्गिक
संख्या~$0$ ही कोणत्याही नैसर्गिक संख्येची उत्तरवर्ती संख्या नाही.
~$f$ ची व्याप्ती सर्व धन पूर्णांकांचा संच $\Int^{+}$ आहे.
\end{ex}

\begin{ex}\ollabel{examplefunext}
$g \colon \Nat \to \Nat$ ची व्याख्या $g(x) = x+2-1$ अशी करू.
यावरून $g$ हे नैसर्गिक संख्या आदान म्हणून घेऊन नैसर्गिक संख्या प्रदान
करणारे फलन आहे हे कळते. नैसर्गिक संख्या x दिली, की $g$ हा~$x$ च्या
उत्तरवर्ती संख्येच्या उत्तरवर्ती संख्येची पूर्ववर्ती संख्या,
म्हणजेच $x+1$, देईल.
\end{ex}
%READERNOTE{OLFUN-003 — गोठवलेल्या इंग्रजी स्रोताच्या गद्यामध्ये याच उदाहरणात आदानासाठी आधी $n$ आणि लगेच x वापरले आहेत. सूत्र न बदलता मराठी गद्याने x हे एकच नाव सातत्याने वापरले आहे.}

\begin{explain}
आपण नुकतीच वेगवेगळ्या \emph{व्याख्या} असलेली $f$ आणि $g$ ही दोन
फलने पाहिली. पण ती \emph{एकच फलन} आहेत. कारण कोणत्याही नैसर्गिक
संख्या~$n$ साठी $f(n) = n+1 = n+2-1 = g(n)$ असते. दुसऱ्या शब्दांत,
$f$ आणि~$g$ यांच्या व्याख्या वेगवेगळ्या समीकरणांनी एकच संगती सांगतात.
म्हणून आपण अप्रत्यक्षपणे फलनांसाठीच्या विस्तारात्मकतेच्या तत्त्वावर
अवलंबून आहोत:
\[
  \text{जर }\forall x\, f(x) = g(x)\text{, तर }f = g
\]
मात्र $f$ आणि~$g$ यांचा प्रांत आणि सहप्रांत एकच असला पाहिजे.
\end{explain}

\begin{ex}
प्रकरणे पाडूनही फलनाची व्याख्या करता येते. उदाहरणार्थ,
$h \colon \Nat \to \Nat$ ची व्याख्या अशी करता येईल:
\[
h(x) =
\begin{cases}
  \frac{x}{2} & \text{जर $x$ सम असेल} \\
  \frac{x+1}{2} & \text{जर $x$ विषम असेल.}
\end{cases}
\]
प्रत्येक नैसर्गिक संख्या सम किंवा विषम असते, म्हणून या फलनाचे प्रदान
नेहमी नैसर्गिक संख्या असेल. प्रकरणे पाडून फलनाची व्याख्या करताना
प्रत्येक शक्य आदान नेमक्या एका प्रकरणात बसले पाहिजे, हे लक्षात ठेवा.
काही वेळा सर्व शक्यता त्या प्रकरणांत येतात आणि कोणतीही शक्यता दोन
प्रकरणांत येत नाही याची सिद्धता करावी लागेल.
\end{ex}

\end{document}
